\subsection{Overview}
We model route generation as a decision--execution hierarchy with a latent skeleton $z$ (e.g., waypoints).
The conditional distribution factorizes as:
\begin{equation}
  p_\theta(\tau \mid c) = \sum_{z} p_\theta(z \mid c)\, p_\theta(\tau \mid z, c).
\end{equation}
In practice, we sample $z$ and then generate $\tau$ conditioned on $z$.

\subsection{Decision: waypoint skeleton}
We represent the route skeleton as an ordered waypoint sequence $z = (w_1,\dots,w_K)$.
An autoregressive planner models $p_\theta(z \mid c)$ to maintain multi-step consistency.

\subsection{Macro guidance: census-derived spatial priors (optional)}
To support semantic controllability without requiring explicit per-agent attribute labels, we allow census-derived spatial covariates (e.g., tract-level population or income) to enter the model as additional context channels or weak guidance terms.
In the main experiments, census guidance is treated as an optional ablation rather than the paper's primary task.

\subsection{Execution: continuous route generation}
Conditioned on $z$, the execution module generates a full-resolution route, modeling local geometry and motion.
We leave the specific choice (diffusion / flow matching) as an implementation detail and focus on the hierarchy and conditioning interfaces.

\subsection{Soft map priors for feasibility}
We incorporate map-derived signals (e.g., road proximity) as soft priors in training and/or sampling, avoiding hard masking of the output space.
This design prevents the generator from collapsing its support to an externally defined feasible set while still discouraging physically implausible routes.
